El volumen acumulado de negociación muestra que
TSLA
fue la empresa con la mayor actividad bursátil durante el período de análisis, acumulando aproximadamente
168.48 mil millones de acciones negociadas
. En segundo lugar se ubicó
AAPL
con
130.58 mil millones
, seguida por
AMD
con
97.65 mil millones
. Estos resultados reflejan que dichas empresas concentraron una parte importante de la actividad de negociación del mercado, evidenciando un elevado interés por parte de los inversionistas y una alta liquidez de sus acciones.
Asimismo, se observa una diferencia considerable entre el primer y el décimo lugar del ranking. Mientras TSLA acumuló cerca de
168.48 mil millones
de acciones negociadas, la décima empresa registró aproximadamente
42.61 mil millones
, lo que representa apenas una cuarta parte del volumen alcanzado por TSLA. Esto evidencia que la actividad de negociación estuvo concentrada en un reducido grupo de compañías de gran capitalización y elevada liquidez.
Es importante señalar que un mayor volumen de negociación no implica necesariamente un mayor rendimiento de la acción. El volumen representa la cantidad de acciones que cambiaron de manos durante el período, por lo que constituye un indicador de liquidez e interés del mercado, pero no garantiza incrementos sostenidos en el precio de la acción.